% GENERATED FILE. Edit canonical lessons, then run npm run generate:books.
% canonical-source-hash: fnv1a64:bb9e377daa117f4c
% canonical-lessons: LA-C353-ferveo, LA-C353-caccabus, LA-C353-elixus, LA-C353-spuma, LA-C353-despumo

\chapter{Aqua fervet}
\label{ch:la-aqua-fervet}
\hlchaptermodality{353}

\begin{chapteropening}
\textbf{By the end of this chapter:} \emph{I can describe boiling food: to boil, the cooking pot, boiled, the foam and to skim.}

Complete the last lesson of chapter 353: I can describe boiling food: to boil, the cooking pot, boiled, the foam and to skim.
\end{chapteropening}

\section[ferveō, fervēre]{ferveō, fervēre — to boil, to be hot}

\label{lesson:LA-C353-ferveo}



\begin{quote}
Before the new one: say the Latin for leaven, then the Latin for a baking pot.
\end{quote}

\subsection*{You'll want to know: ferveō, fervēre}
\textbf{ferveō, fervēre} — \textquotedblleft{}to boil, to be hot\textquotedblright{}.

\begin{quote}
Aqua in ōllā fervēbat.
\end{quote}

\begin{quote}
\emph{The water was boiling in the pot.}
\end{quote}

\emph{ferveō, fervēre, ferbuī}, also an older \emph{fervō, fervere}. English \emph{fervent}.

\begin{cousinweb}
\textbf{ferveō} is an inherited verb (EtymDB). English \emph{fervent} comes from it, and \emph{ferment} from its noun \emph{fermentum}.
\end{cousinweb}

\begin{grammarlens}[title={What you've built}]
One more word for telling, explaining and describing at B1.
\end{grammarlens}

\subsection*{Your turn}
\begin{itemize}
\raggedright
  \item \emph{Say it:} \emph{ferveō, fervēre}
  \item \emph{Say it:} \emph{ferveō, fervēre}, once more
  \item \emph{Say it:} a sentence of your own with \emph{ferveō, fervēre}
  \item \emph{Recall:} say the Latin for leaven, then the Latin for a baking pot, then say \emph{ferveō, fervēre} again
\end{itemize}

\begin{tcolorbox}[breakable,colback=teal!4,colframe=teal!35!black,title={Before you move on}]
What does ferveō, fervēre mean? (\textquotedblleft{}To boil, to be hot\textquotedblright{}.) Say it once more.
\end{tcolorbox}
\section[caccabus, caccabī]{caccabus, caccabī — a cooking pot}

\label{lesson:LA-C353-caccabus}



\begin{quote}
Before the new one: say the Latin for a baking pot, then the Latin for to boil.
\end{quote}

\subsection*{You'll want to know: caccabus, caccabī}
\textbf{caccabus, caccabī} — \textquotedblleft{}a cooking pot\textquotedblright{}.

\begin{quote}
Caccabus super focum pendēbat.
\end{quote}

\begin{quote}
\emph{A cooking pot was hanging over the hearth.}
\end{quote}

Masculine, a Greek loan: a round pot of bronze or clay.

\begin{cousinweb}
\textbf{caccabus} was borrowed from Greek \emph{kakkabos}, \textquotedblleft{}a pot, a pan\textquotedblright{} (EtymDB).
\end{cousinweb}

\begin{grammarlens}[title={What you've built}]
One more word for telling, explaining and describing at B1.
\end{grammarlens}

\subsection*{Your turn}
\begin{itemize}
\raggedright
  \item \emph{Say it:} \emph{caccabus, caccabī}
  \item \emph{Say it:} \emph{caccabus, caccabī}, once more
  \item \emph{Say it:} a sentence of your own with \emph{caccabus, caccabī}
  \item \emph{Recall:} say the Latin for a baking pot, then the Latin for to boil, then say \emph{caccabus, caccabī} again
\end{itemize}

\begin{tcolorbox}[breakable,colback=teal!4,colframe=teal!35!black,title={Before you move on}]
What does caccabus, caccabī mean? (\textquotedblleft{}A cooking pot\textquotedblright{}.) Say it once more.
\end{tcolorbox}
\section[ēlixus, ēlixa, ēlixum]{ēlixus, ēlixa, ēlixum — boiled}

\label{lesson:LA-C353-elixus}



\begin{quote}
Before the new one: say the Latin for to boil, then the Latin for a cooking pot.
\end{quote}

\subsection*{You'll want to know: ēlixus, ēlixa, ēlixum}
\textbf{ēlixus, ēlixa, ēlixum} — \textquotedblleft{}boiled\textquotedblright{}.

\begin{quote}
Carō ēlixa mollis et dulcis erat.
\end{quote}

\begin{quote}
\emph{The boiled meat was soft and sweet.}
\end{quote}

Boiled right through. A slave in Plautus jokes that he tastes better boiled than roasted.

\begin{cousinweb}
\textbf{ēlixus} is built, says Lewis \& Short, on \emph{ē}, \textquotedblleft{}out, thoroughly\textquotedblright{}, and \emph{lix}, an old word for water or lye: thoroughly boiled.
\end{cousinweb}

\begin{grammarlens}[title={What you've built}]
One more word for telling, explaining and describing at B1.
\end{grammarlens}

\subsection*{Your turn}
\begin{itemize}
\raggedright
  \item \emph{Say it:} \emph{ēlixus, ēlixa, ēlixum}
  \item \emph{Say it:} \emph{ēlixus, ēlixa, ēlixum}, once more
  \item \emph{Say it:} a sentence of your own with \emph{ēlixus, ēlixa, ēlixum}
  \item \emph{Recall:} say the Latin for to boil, then the Latin for a cooking pot, then say \emph{ēlixus, ēlixa, ēlixum} again
\end{itemize}

\begin{tcolorbox}[breakable,colback=teal!4,colframe=teal!35!black,title={Before you move on}]
What does ēlixus, ēlixa, ēlixum mean? (\textquotedblleft{}Boiled\textquotedblright{}.) Say it once more.
\end{tcolorbox}
\section[spūma, spūmae]{spūma, spūmae — foam, froth, scum}

\label{lesson:LA-C353-spuma}



\begin{quote}
Before the new one: say the Latin for a cooking pot, then the Latin for boiled.
\end{quote}

\subsection*{You'll want to know: spūma, spūmae}
\textbf{spūma, spūmae} — \textquotedblleft{}foam, froth, scum\textquotedblright{}.

\begin{quote}
Spūma alba in aquā ferventī natābat.
\end{quote}

\begin{quote}
\emph{White foam was floating on the boiling water.}
\end{quote}

Feminine. English \emph{spume}.

\begin{cousinweb}
\textbf{spūma} is linked by Lewis \& Short with \emph{spuō}, \textquotedblleft{}to spit\textquotedblright{}. English \emph{spume} comes from it.
\end{cousinweb}

\begin{grammarlens}[title={What you've built}]
One more word for telling, explaining and describing at B1.
\end{grammarlens}

\subsection*{Your turn}
\begin{itemize}
\raggedright
  \item \emph{Say it:} \emph{spūma, spūmae}
  \item \emph{Say it:} \emph{spūma, spūmae}, once more
  \item \emph{Say it:} a sentence of your own with \emph{spūma, spūmae}
  \item \emph{Recall:} say the Latin for a cooking pot, then the Latin for boiled, then say \emph{spūma, spūmae} again
\end{itemize}

\begin{tcolorbox}[breakable,colback=teal!4,colframe=teal!35!black,title={Before you move on}]
What does spūma, spūmae mean? (\textquotedblleft{}Foam, froth, scum\textquotedblright{}.) Say it once more.
\end{tcolorbox}
\section[dēspūmō, dēspūmāre]{dēspūmō, dēspūmāre — to skim off}

\label{lesson:LA-C353-despumo}



\begin{quote}
Before the new one: say the Latin for boiled, then the Latin for foam.
\end{quote}

\subsection*{You'll want to know: dēspūmō, dēspūmāre}
\textbf{dēspūmō, dēspūmāre} — \textquotedblleft{}to skim off\textquotedblright{}.

\begin{quote}
Coquus ōllam dīligenter dēspūmāvit.
\end{quote}

\begin{quote}
\emph{The cook carefully skimmed the pot.}
\end{quote}

\emph{dēspūmō, dēspūmāre, dēspūmāvī, dēspūmātum}. Vergil's farmer skims the boiling grape juice with leaves.

\begin{cousinweb}
\textbf{dēspūmō} is built on \emph{spūma}, \textquotedblleft{}foam\textquotedblright{}, with \emph{dē-}, \textquotedblleft{}off\textquotedblright{} (EtymDB): to take the foam off.
\end{cousinweb}

\begin{grammarlens}[title={What you've built}]
One more word for telling, explaining and describing at B1.
\end{grammarlens}

\subsection*{Your turn}
\begin{itemize}
\raggedright
  \item \emph{Say it:} \emph{dēspūmō, dēspūmāre}
  \item \emph{Say it:} \emph{dēspūmō, dēspūmāre}, once more
  \item \emph{Say it:} a sentence of your own with \emph{dēspūmō, dēspūmāre}
  \item \emph{Recall:} say the Latin for boiled, then the Latin for foam, then say \emph{dēspūmō, dēspūmāre} again
\end{itemize}

\begin{tcolorbox}[breakable,colback=teal!4,colframe=teal!35!black,title={Before you move on}]
What does dēspūmō, dēspūmāre mean? (\textquotedblleft{}To skim off\textquotedblright{}.) Say it once more.
\end{tcolorbox}
